% !TeX root

\documentclass[12pt,a4paper]{article}
% =============================================================
% Preamble entry point. Each file in config/ handles one concern.
% Order matters: do not reorder without checking package loading.
% =============================================================
% FONTS & LANGUAGE
% Font encoding: T1 for proper accents and hyphenation
\usepackage[T1]{fontenc}
% Language: British English (matches A4 paper).
% For American English, change to main=american.
\usepackage[main=british]{babel}
% Fonts: Times-like text and math (newtx)
\usepackage{newtxtext,newtxmath}
% Quotes: style follows the babel language; write \enquote{...} in the text.
% Must be loaded after babel and before biblatex.
\usepackage[autostyle=true]{csquotes}
% Typography: better justification and hyphenation
\usepackage{microtype}

% Note: no inputenc needed, UTF-8 is the default in modern LaTeX.
          % fonts, babel (English), csquotes
% PAGE GEOMETRY & TITLE POSITION
% 1-inch margins on all sides
\usepackage[a4paper,margin=1in]{geometry}
% Title block hooks (\pretitle, \posttitle, \predate, ...); load before authblk
\usepackage{titling}
\setlength{\droptitle}{-5em}

% COLORS & TEXT LAYOUT
% Named SVG colors and table row coloring
\usepackage[table, svgnames]{xcolor}
% Line spacing: \onehalfspacing, \doublespacing
\usepackage{setspace}
% Paragraphs: vertical space instead of indentation
\usepackage{parskip}
% Comment blocks: \begin{comment} ... \end{comment}
\usepackage{comment}

% LISTS
\usepackage{enumitem}
% Tighter itemize spacing
\setlist[itemize]{topsep=0pt, partopsep=0pt, parsep=0pt}

% HEADERS & FOOTERS
\usepackage{fancyhdr}
\setlength{\headheight}{14.49998pt} % Fixes fancyhdr warning with 12pt text
\fancyhf{}
% Footer: "X of Y", centered. The label mylastpage is set in config/07_macros.tex
\fancyfoot[C]{\thepage\ of \pageref{mylastpage}}
\pagestyle{fancy}
% Plain pages (title page, chapter starts): no header, no rules
\fancypagestyle{plain}{
  \fancyhf{}
  \fancyfoot[C]{\thepage\ of \pageref{mylastpage}}
  \renewcommand{\headrulewidth}{0pt}
  \renewcommand{\footrulewidth}{0pt}
}

% TITLE BLOCK
% Multiple authors via \author{...} and \affil{...}
\usepackage{authblk}
\renewcommand\Authfont{\small}
% Separate authors with commas instead of "and"
\renewcommand\Authsep{, }
\renewcommand\Authand{, }
\renewcommand\Authands{, }
            % colors, spacing, headers/footers, lists, title block
% Math environments (align, gather, ...) and symbols
\usepackage{amsmath}
% Bold math: \bm{...}
\usepackage{bm}
              % math packages
% GRAPHICS & CAPTIONS
% [H] placement specifier for exact float position
\usepackage{float}
\usepackage{graphicx}
% Keeps figures inside the margins (adjustbox)
\usepackage{adjustbox}
\graphicspath{{img/}} % Folder with image assets
% TODO: remove if no section uses it
\newcounter{figcounter}
% Captions: small font
\usepackage{caption}
\captionsetup{font=small}
% Side-by-side figures with (a), (b) sub-captions
\usepackage{subcaption}

% TIKZ & PLOTS
\usepackage{tikz}
\usepackage{pgfplots}
\pgfplotsset{compat=1.18}
\usetikzlibrary{3d}

% TABLES & CODE
% Code listings
\usepackage{listings}
% tabularray tables
\usepackage{tabularray}
% Tabularray templates: no repeated head/foot across page breaks,
% caption placed in the foot
\DefTblrTemplate{firsthead}{default}{}
\DefTblrTemplate{middlehead}{default}{}
\DefTblrTemplate{lasthead}{default}{}
\DefTblrTemplate{firstfoot}{default}{\UseTblrTemplate{caption}{default}}
\DefTblrTemplate{foot}{default}{\UseTblrTemplate{caption}{default}}
    % graphics, captions, TikZ, tables, listings
% Requires babel and csquotes to be loaded first (config/01_language.tex).
\usepackage[backend=biber, style=alphabetic, sorting=ydnt]{biblatex}
\addbibresource{references.bib} % Bibliography database
      % biblatex/biber
% Keep this file late in setup.tex. Packages that need hyperref (e.g. cleveref) must be loaded after it.

% URL line breaking (xurl also loads url)
\usepackage{xurl}
% Clickable references, citations and URLs
\usepackage[colorlinks=true, linkcolor=blue, citecolor=blue, urlcolor=blue, pdfborder={0 0 0}, hyperfootnotes=true]{hyperref}
% Links to figures/tables jump to the top of the float (needs hyperref)
\usepackage[table, figure]{hypcap}
        % url + hyperref (keep late)
% CUSTOM COMMANDS

% Clickable monospaced email: \emailaddr{name@example.com}
\newcommand{\emailaddr}[1]{\href{mailto:#1}{\texttt{#1}}}

% End of document with bibliography.
% The last-page label comes after the bibliography, so "X of Y" counts it.
\newcommand{\conclusion}{
  \clearpage
  \printbibliography
  \label{mylastpage}
}

% End of document with bibliography kept outside the page count.
% The label comes before the bibliography, and the bibliography pages have no footer.
\newcommand{\conclusionseparate}{
  \label{mylastpage}
  \clearpage
  \pagestyle{empty}
  \printbibliography
}

% End of document without bibliography (documents with no citations).
% Keeps the "X of Y" footer working.
\newcommand{\conclusionnobib}{
  \label{mylastpage}
}

% Compact title block for documents with no authors and no date:
% removes the empty author/date lines and pulls the text up under the title.
% Use in the preamble right after \title{...}
\newcommand{\compacttitle}{
  \author{}\date{}
  \predate{}\postdate{}
  \posttitle{\par\end{center}\vspace{-2em}}
}
            % custom commands


\title{The title of your report}

\author{names of students}
\affil{} % Necessary

\date{\today}

\begin{document}

\maketitle
\tableofcontents

\section{Introduction} \label{sec:intro}

This section is to give a general introduction to the company, including a description of its main \textit{product offerings},
\textit{markets} and \textit{strategy}.

\section{Value Network Analysis} \label{sec:value}

This section is to show the \textbf{value} network in which the company will work.
\textit{Agents}, \textit{resources}, and \textit{exchanges of resources} are to be made explicit.
Note that the project group here has to clarify exactly which products the company is to offer.
Also note that the value network analysis has to take into account that the company is to work in different markets with different
actors.

\subsection{VDML Diagram}

Include a Value Delivery Modeling Language (\cite{OMG_VDML}) diagram here.

\subsection{VDML Diagram Explained}

Include a textual description explaining the VDML diagram.

\section{Goal Design} \label{sec:goals}

This section is to show the \textbf{goals} of the company.
The goals shall be clearly related to the value network analysis in \autoref{sec:intro}.
In order to identify relevant goals, the project group shall make use of the five forces model \cite{QuickMBA_FiveForces} and generic
strategies of Michael Porter (look on the web).
It is recommended that Strategy maps be used to structure the goal model \cite[ch.~2, pp.~29--52]{Kaplan2004_StrategyMaps}.

\subsection{Goal Model}

Include a graphical Business Motivation Model (\cite{BusinessRulesGroup_BMM}) goal model here.

\subsection{Goal Model Explained}

Include a textual description of the BMM model that explains it and shows how it is related to the components of the five forces
model.

\section{Process Design} \label{sec:processes}

This section is to show the \textbf{processes} of the company.
The processes shall be clearly related to the value network analysis in \autoref{sec:value} and the goal design in
\autoref{sec:goals}.

\subsection{Exchange and Conversion Processes in REA}

Include a graphical conceptual schema of the company with exchange and conversion processes (expressed as UML class diagrams)
using the REA (Resources, Events and Agents) ontology.
All classes including attributes must be named and the cardinality/multiplicity between the classes must be expressed.
Stereotypes for all classes shall be specified in accordance with REA.
Note that only the exchange and conversion processes within the company shall be modelled, but not processes in other companies in
the value network.
The conceptual schema shall also make use of other relevant patterns than those for exchange and conversion processes, i.e.
structural and behavioural patterns as defined in the book by Hruby \cite{Hruby2006_BusinessPatterns}.
Provide textual descriptions for explaining the graphical schema.

\subsection{Detailed Processes in EPC}

Include a number of detailed process models, expressed as EPC (Event-driven Process Chains) diagrams.
The processes should include sales and procurement.
The process models should show actions of individual actors and handle alternatives and exceptions.
It is recommended to make each process diagram small by using the decomposition mechanisms of EPC diagrams.

\subsection{Relationships between Models}

Include one or several tables that visualize the relationships between the models introduced in this section and the models in the
previous sections.
The tables should help to answer the following questions:
\begin{itemize}
  \item How does the design of the processes help in fulfilling the goals of the company?
  \item What are the relationships between the exchange and conversion processes and the EPC models?
  \item What are the relationships between the EPC models and the VDML diagram?
\end{itemize}

\section{IT Architecture Design} \label{sec:it_architecture}

This section is to show the IT architecture of the company.
The IT systems to be used in the company are to be specified, including BPMS, CRM, ERP, DW, BI tools, dashboard, and KM.
The use of EAI (Enterprise Application Integration) and decision support systems (BI, DW and Dashboards) shall also be discussed.
The \textbf{IT architecture} shall be documented graphically in a \textbf{diagram} showing the systems and the data flows between
them, representing systems as nodes and data flows as labeled arrows.
The relationships between the systems and the business processes they support shall be made explicit.

\section{Business Performance Management} \label{sec:performance}

This section is to introduce the key performance indicators (\textbf{KPI}) for measuring and analysing the business processes of
the company.

\subsection{Leading and Lagging KPIs}

Include between five and ten leading and lagging indicators.
Each indicator shall be described in an indicator template consisting of at least the following parts: name of the indicator,
definition of the indicator, which goal(s) in the company's goal model that the indicator is related to, target values, which
part of a business process (i.e. which activity) the indicator measures, which IT system that will be the data source for the
indicator, and which type of IT solution will feed the dashboard system with data from the source system.
Ensure that the descriptions are consistent with the IT architecture in \autoref{sec:it_architecture}.

\subsection{Evaluation of KPIs}

Include an evaluation of the KPIs from the previous subsection, based on the article by Eckerson
\cite[ch.~11, pp.~197--216]{Eckerson2010_Performance}.

\section{Enterprise Architecture Evaluation} \label{sec:vulnerabilities}

In other parts of the work you have designed the organization from several perspectives, goals, processes, IT architecture, and so
on.
All of these combined converge towards an Enterprise Architecture (EA).
A main task of EA is to provide an abstracted overview of the organization to support the management in their decision making.
Therefore, it is necessary to have an \textbf{EA model} at hand to be able to answer upcoming questions.
In this light, your task is to create an EA model that includes the key parts of the previously created models.
Security becomes more and more an important topic for organizations due to the changing political situation in the world and the
bigger exposure of infrastructure to the internet (Industrial Internet of Things (IIoT)).
Therefore, the management asks you for an assessment of how the goals of the organization are threatened by actual security
issues.
To answer this question, analyze reported vulnerabilities within your IT Architecture (e.g., by using the National Vulnerability
Database (NVD)) and determine their influence on the goals by analyzing the connection within your EA model
\cite{Ebbers2021_IIoTVulnerabilities}.

\conclusion

\end{document}